\chapter{Conclusion and Future Scope}

Financial Fraud Detection Using Graph Machine Learning demonstrates that treating a transaction history as a network rather than as a list of independent events changes what can be detected. The system successfully achieves the objectives set at the outset. Transactions are modelled as a directed weighted graph in which accounts are nodes, and NetworkX computes PageRank, degree centrality, Louvain community membership and cycle membership for every account. A transparent rule engine applies six checks derived from standard anti-money-laundering typologies, each alert carrying the rule code and a readable reason. A feature engineering module combines eleven transaction-level aggregates with six graph-derived metrics into a seventeen-feature vector per account, and Random Forest, XGBoost and Isolation Forest models score every account from that matrix. Graph Convolutional Network and GraphSAGE models learn from the account graph directly alongside them. The imbalance of the fraudulent class is addressed by a different mechanism in each family --- balanced class weights for Random Forest, a positive-class weight for XGBoost, a contamination rate for the unsupervised Isolation Forest, and a weighted cross-entropy loss for the two graph neural networks. Detection output is carried through to a complete investigation workflow in which alerts are grouped into fraud rings with member roles assigned, an analyst confirms or clears each account against an audit trail, the holder of a confirmed account is notified by email, a further transaction sent by a confirmed account pushes a real-time alert to the analyst --- a notification, not a block on the payment --- and cases are exported as PDF reports. A React single-page application provides the analyst interface over a FastAPI service and a SQLite database accessed entirely through the SQLAlchemy ORM.

The system comprises a FastAPI backend of eighteen routers and seven detection modules, a React client of eighteen pages, an automated test suite of fifty-nine tests across unit, integration and workflow layers, and a public deployment with the backend and frontend hosted separately. Evaluated on the seeded synthetic dataset, the rule engine reached 96.21\% recall at 99.88\% precision, GraphSAGE reached an F1 score of 0.839 against XGBoost's 0.825, and 52 fraud rings were assembled with 96.8\% of their members genuinely fraudulent. On the independent Elliptic Bitcoin benchmark GraphSAGE reached an F1 score of 0.549. These figures are drawn from the evaluation outputs the application itself writes, and the caveats attached to them in Chapter~\ref{ch:results} --- a synthetic primary dataset, two different test partitions, and evaluations that hold out labels rather than accounts, so that no result here speaks to performance on an account the graph had never seen --- apply equally here. A verification re-run from the committed code and data reproduced the rule-engine, Random Forest, Isolation Forest and fraud-ring results exactly and the XGBoost result to within 0.006 in F1 (0.831 on re-run); the graph neural network figures rest on stored result files that could not be regenerated in that verification. The roughly twenty-two-second pipeline runtime is an informal observation from the development laptop rather than a benchmark; the verification re-run, on different hardware, completed in 17 to 19 seconds.

In terms of future enhancements, the most immediate priority is evaluation on a real labelled transaction graph rather than a synthetic one. The gap between the synthetic and the Elliptic results reported in Chapter~\ref{ch:results} shows that performance on planted structures says little about performance in the field, and only real data settles it. Closely related is defending the graph neural networks against camouflage: both currently weight every neighbour alike, so a fraudster who transacts widely with legitimate accounts dilutes their own signal, and a neighbour selection mechanism of the kind CARE-GNN introduced would address it. Incremental graph updates --- recomputing metrics for the affected neighbourhood rather than the whole graph --- would remove the principal scalability constraint and are a prerequisite for the streaming detection path that would allow a transaction to be scored as it arrives rather than after the fact. On the engineering side, pinning dependency versions, archiving each evaluation's result files with the commit that produced them, and committing the training-only Isolation Forest contamination estimate would make every figure in this report reproducible from the repository alone.

In the medium term, multi-tenant support through institution-namespaced data, an analyst feedback loop in which case resolutions are fed back as labels for model retraining, and configurable rule thresholds exposed through the interface rather than fixed in code would all move the system closer to operational use. Further ahead, temporal graph modelling that treats the transaction network as evolving over time rather than as a single static snapshot, and confidence scoring for ring assembly, would address the two most significant methodological limitations identified in Chapter~\ref{ch:results}.

The foundational architecture --- a graph representation of the transaction history, a layered detection pipeline with independently runnable stages, an ORM-mediated persistence layer, and a separation between detection and investigation --- is well suited to support all of these enhancements without requiring a re-architecture. Future work will focus on progressively addressing the limitations identified in Chapter~\ref{ch:results}, beginning with evaluation on real labelled data and camouflage-resistant graph learning.
